% RESEARCH_DRAFT_SHA256: 333dc19c6ed2bca67f7115ad5aaef257c159d7dcb35e954d7b24ecd524d54257
% Layout and source-reference wrappers expanded for the web renderer.
\documentclass[11pt]{article}
\usepackage[margin=25mm]{geometry}
\usepackage{fontspec,unicode-math}
\setmainfont{Latin Modern Roman}
\setmathfont{Latin Modern Math}
\usepackage{amsmath,amssymb,mathtools,xurl,hyperref}
\hypersetup{colorlinks=true,urlcolor=blue,linkcolor=blue}
\setlength{\parindent}{0pt}
\setlength{\parskip}{5pt}
\setlength{\emergencystretch}{3em}
\newcommand{\sref}[2]{\hyperlink{source:#1:#2}{[S#2]}}
\newcommand{\cataloguescope}[1]{\paragraph{Source question.}#1}
\newcommand{\floor}[1]{\left\lfloor#1\right\rfloor}
\begin{document}

\section*{Mixed-integer bipartite optimization: exact easy hulls and a mixing obstruction}
\noindent catalogue entry; catalogue entry. Research attempt dated 11 October 2026.
\paragraph{Result in scope.}This manuscript contains a source-specific attempt and complete proofs of the stated auxiliary calculations. It makes no new solution or novelty claim for the entire target.
\section{Original problem and preserved definition}
The following definition and source list are reproduced from the repository. Any earlier inspection dates and reported literature claims in this preserved text describe that source record, not additional verification by this attempt.
% BEGIN PRESERVED DEFINITION
\subsection{Definitions and mathematical statement}

\paragraph{Shared notation.}Standard number systems are denoted by $\mathbb N,\mathbb Z,\mathbb Q,\mathbb R,\mathbb C$; index ranges, computational models and specialized notation follow the cited problem source and the definitions below.


A linear program optimizes $c^Tx$ over a polyhedron $\{x:Ax\leq b\}$. In the simplex method a pivot moves between adjacent basic feasible solutions; a pivot rule chooses the next such move.

An integer program imposes integral coordinates on a system of constraints. A mixed-integer program allows both integral and real coordinates; $\operatorname{conv}(S)$ denotes the set of finite convex combinations of feasible points of $S$.
\paragraph{Statement recorded in the supplied source.}
For a bipartite graph $G=(U\cup V,E)$, an integral-coordinate set $I$, and $S(G,I)=\{x\in\mathbb{R}^{U\cup V}:x_i+x_j\geq b_{ij}\text{ for }ij\in E,\ x_i\in\mathbb{Z}\text{ for }i\in I\}$, determine the complexity of optimizing a linear function over $S(G,I)$, equivalently deciding whether a given $\bar{x}$ belongs to $\operatorname{conv}(S(G,I))$. \href{https://doi.org/10.4171/owr/2008/51}{[S2]}
\paragraph{Scope qualification.}
The report\textquotesingle s encoding and assumptions on $b_{ij}$ and boundedness are retained; the asserted optimization/membership equivalence is not independently established here.
\paragraph{Source status.}
The supplied catalogue labels this record open. Its source-status metadata was checked against the supplied snapshot on 2026-10-11. This import does not independently establish current open status or mathematical correctness. The catalogue literature assessment is dated 2026-08-21.
\subsection{Short English statement}
Classify linear optimization and convex-hull membership for mixed-integer covering constraints on bipartite graphs.
\subsection{Sources}
\begin{itemize}
\item[S1] Ulam AI and the UnsolvedMath Contributors, UnsolvedMath: A Curated Collection of Open Mathematics Problems, catalogue entry. Catalogue curation and metadata: CC BY 4.0. The supplied source question is retained; definitions and the English summary are editorial additions.
\url{https://huggingface.co/datasets/ulamai/UnsolvedMath}.
\item[S2] Combinatorial Optimization (2009), Oberwolfach Reports, DOI 10.4171/owr/2008/51, pp. 2875--2942
\url{https://doi.org/10.4171/owr/2008/51}.
\item[S3] network formulations
\url{https://doi.org/10.1287/moor.1080.0354}.
\item[S4] mixed-integer bipartite covers
\url{https://doi.org/10.1287/moor.1100.0454}.
\textit{Underlying Oberwolfach Reports material retains its source-specific licence, including CC BY-SA 4.0 where applicable; see THIRD\_PARTY\_NOTICES.md.}
\end{itemize}
% END PRESERVED DEFINITION
\subsection{Verbatim supplied catalogue question}
For a bipartite graph $G=(U\cup V,E)$, an integral-coordinate set $I$, and $S(G,I)=\{x\in\mathbb{R}^{U\cup V}:x_i+x_j\geq b_{ij}\text{ for }ij\in E,\ x_i\in\mathbb{Z}\text{ for }i\in I\}$, determine the complexity of optimizing a linear function over $S(G,I)$, equivalently deciding whether a given $\bar{x}$ belongs to $\operatorname{conv}(S(G,I))$.

\section{Attempted reduction to ordinary linear programming}
The source's $b_{ij}$ are rational, $k$ is their least common denominator,
and variables have no nonnegativity constraints unless separately stated.
We try to recover an exact ordinary linear relaxation by bipartite rounding.
It succeeds in two boundary regimes and fails on a small mixed example.

\paragraph{Integral right-hand sides and fully integral variable sets.}
For integer $b_{uv}$, any relaxed feasible $x$ can be decomposed as a
finite convex combination of fully integral feasible points. To prove it,
put $y_u=x_u$ on $U$ and $y_v=-x_v$ on $V$, so
$y_u-y_v\ge b_{uv}$. Round all $y_i$ to
$z_i(\theta)=\lfloor y_i+\theta\rfloor$ using one uniform
$\theta\in[0,1)$. Integer translation of floors implies
$z_u-z_v\ge b_{uv}$ on every edge. Undoing the sign change gives an
integer feasible point for every $\theta$. For $y_i=m+r$ with
$m\in\mathbb Z$ and $0\le r<1$, the rounded value is $m$ for a
set of parameters of length $1-r$ and $m+1$ for length $r$, hence
has mean $y_i$. Partitioning $[0,1)$ at the finitely many coordinate
thresholds makes this a finite convex combination with mean $x$.
Thus for every $I$ and integer right-hand sides,
\[
\operatorname{conv}(S(G,I))=\{x:x_u+x_v\ge b_{uv}\ \forall uv\in E\}.
\]
If instead $I=U\cup V$ and the right-hand sides are rational, replace
every $b_{uv}$ by $\lceil b_{uv}\rceil$: an integer sum satisfies the
original inequality exactly when it satisfies the rounded one. The preceding
proof then gives the hull for that regime as well. If $I$ is empty, the
original set is already convex. Membership in these displayed hulls is
checked by their edge inequalities. Optimization reduces to ordinary rational
linear programming, with its standard polynomial-time algorithms as an
external computational dependency; no new LP algorithm is claimed.

\paragraph{A three-vertex mixed example blocks the same relaxation.}
Use a path with integral endpoints $x_1,x_2\in\mathbb Z$ and continuous
center $z\in\mathbb R$, with
\[
x_1+z\ge1/3,\qquad x_2+z\ge2/3.
\]
The denominator parameter is three. Every mixed feasible point satisfies
\[
x_1+x_2+2z\ge4/3.
\]
Indeed put $m=x_2-x_1\in\mathbb Z$ and $w=x_1+z$. The constraints
say $w\ge\max(1/3,2/3-m)$, so the objective is $m+2w$.
If $m\le0$, this is at least $4/3-m\ge4/3$.
If $m\ge1$, it is at least $m+2/3\ge5/3$.
Equality $4/3$ is attained by $x_1=x_2=0,z=2/3$.
In contrast $(x_1,x_2,z)=(1/3,2/3,0)$ is feasible for the ordinary
continuous relaxation and has objective one. Since the stronger inequality
is preserved by convex combinations, this point is outside the mixed hull.
The objective is bounded below despite unrestricted variable signs, as the
calculation shows; no unmentioned nonnegativity assumption creates the gap.

\paragraph{Denominator enumeration is not a bit-polynomial conclusion.}
The report studies extended formulations whose size grows with $k$.
Even one right-hand side $b=1/2^L$ has a binary encoding of length $O(L)$
while its least denominator is $k=2^L$. Therefore an algorithm that
explicitly creates $k$ states or variables may take exponential time in
the original binary input length. Calling it polynomial in $k$ does not
classify the source's general encoded-input problem. This numerical example
is exact and requires no asymptotic number-theoretic assumption.

\section{Adversarial audit and unresolved complexity}
The audit checked unrestricted-sign variables, use of the common random
threshold, finite rather than merely limiting mixtures, and the denominator
encoding. The supplied optimization/membership equivalence requires a precise
oracle/encoding formulation; it is not proved by the edge-wise argument.
The small path demonstrates why a plain-LP proposal fails, but does not show
NP-hardness or prevent another polynomial method. The general rational mixed
case and its complexity remain unresolved by this attempt. No numerical
experiment was performed.

\section{Process and contributions}
Mathematics and model-assisted adversarial self-audit: OpenAI gpt-6-astra, ultra reasoning, as dispatched for this research batch; exact serving revision not exposed. Codex (GPT-6; exact serving revision not exposed) materially assisted with orchestration and format review. No human mathematical authorship, mathematical edits, or independent human proof checking is asserted.
Contribution prompt A was applied to the full definition and supplied source row, then prompt B was applied adversarially to the derived results. The item-specific audit above records the scope checks and rejected extrapolations. Local assembly and structural checks do not establish mathematical correctness. The source catalogue status is preserved in TOP\_PROBLEM and is separate from this attempt status.
\paragraph{Integration review.} Codex (GPT-6; exact serving revision not exposed) reviewed the written arguments and expanded layout and source-reference wrappers. Repository metadata and the production web renderer were checked using Python 3.13 and JavaScript with MathJax 3.2.2. This model-assisted review is not independent human proof verification; no PDF compilation or formal proof checking is asserted.

\section{Completion Estimate}
\noindent\textbf{COMPLETION: 0\%}
The proved material is a rederivation of known or auxiliary facts rather than a substantiated new advance on the original target. The general rational mixed case has nontrivial mixing inequalities and denominator-sensitive formulations; no polynomial-time algorithm, hardness classification, or full optimization/membership reduction is supplied.
\section{References}
The complete original Sources subsection above is retained, including catalogue attribution and licence notices. The additional list records primary sources actually consulted for this attempt on 11 October 2026; uncited source links reproduced above were not all rechecked.
\begin{itemize}
\item Oberwolfach Report 51/2008: Conforti, Combinatorial Mixed-Integer Programming, printed pp.2904--2905; Alcuin-number contribution, pp.2891--2893. Each relevant original passage read from cached primary PDF text on 11 October 2026. \url{https://ems.press/content/serial-article-files/46195}.
\end{itemize}
\section{Final status}
\textbf{Unresolved in this attempt.} The general rational mixed case has nontrivial mixing inequalities and denominator-sensitive formulations; no polynomial-time algorithm, hardness classification, or full optimization/membership reduction is supplied.
\end{document}
